\section{Введение}
\label{sec:Chapter0} \index{Chapter0}

% \textcolor{red}{В этой части надо описать предметную область, задачу из которой вы будете
% решать, объяснить её актуальность (почему надо что-то делать сейчас?). Здесь же
% стоит ввести определения понятий, которые вам понадобятся в постановке задачи.}

Развитие информационной инфраструктуры, начиная от облачных платформ и центров обработки данных,
заканчивая персональными компьютерами, всё чаще полагается на технологии виртуализации.

Виртуализация --- группа технологий, основанных на преобразовании формата или параметров 
программных или сетевых запросов к компьютерным ресурсам с целью обеспечения независимости 
процессов обработки информации от программной или аппаратной платформы информационной системы.
\footnote{
    ГОСТ Р 56938-2016. Национальный стандарт Российской Федерации. Защита информации. 
    Защита информации при использовании технологий виртуализации. Общие положения
}

Простые пользователи используют т.н. виртуальные машины для запуска альтернативных операционных
систем на своих рабочих станциях для тестирования подозрительных программ, для разработки и 
обучения. Облачные провайдеры используют виртуальные машины для обеспечения работы серверных 
платформ: в частности, для предоставления пользователям возможности аренды вычислительных ресурсов.
Последним это позволяет не покупать собственное оборудование для разворачивания инфраструктуры, 
масштабировать мощности под нагрузку и платить только за фактическое использование ресурсов.

Виртуальная машина --- виртуальная вычислительная система, которая состоит из виртуальных устройств 
обработки, хранения и передачи данных и которая дополнительно может содержать программное обеспечение 
и пользовательские данные.
\footnote{
    ГОСТ Р 56938-2016. Национальный стандарт Российской Федерации. Защита информации. 
    Защита информации при использовании технологий виртуализации. Общие положения
}

Отдельным фактором, влияющим на рост спроса на технологии виртуализации, является развитие машинного
обучения и искусственного интеллекта. Развертывание и обучение моделей требует больших вычислительных 
ресурсов (мощных GPU и NPU). Облачные провайдеры предоставляют виртуальные машины с таким 
оборудованием, позволяя бизнесу и частным пользователям проводить исследования, не приобретая
дорогостоящее оборудование. Это дополнительно увеличивает потребность в эффективном использовании 
аппаратного обеспечения, что также достигается при помощи технологий виртуализации.

Развитие технологий виртуализации предъявляет повышенные требования 
к инструментам диагностики и отладки, которые должны обеспечивать сбор
данных из гостевых операционных систем с минимальным влиянием на их
временные характеристики. В задачах логгирования и профилирования ключевыми
являются низкая задержка и высокая пропускная способность передачи
диагностической информации. В связи с этим возникает необходимость
создания специализированного канала связи «гость – хост», оптимизированного
для однонаправленной потоковой передачи данных.

Ещё одной особенностью современных виртуальных машин является наличие 
сопутствующего программного обеспечения, позволяющего пользователям бесшовно 
с ними взаимодействовать. Например, инструменты, обеспечивающие общий буфер 
обмена, перенаправление USB-устройств, синхронизацию времени и разрешения 
экрана. Эти системы в процессе своей работы также генерируют диагностические
события, которые могут быть полезны для выявления ошибок у пользователей.
Таким образом, система сбора диагностической информации должна эффективно
работать с учётом функционирования нескольких самостоятельных компонентов внутри 
виртуальной машины. 

Разработка подобной системы передачи диагностических данных подразумевает
создание драйвера в гостевой операционной системе, который взаимодействует
со специальным виртуальным устройством на стороне эмулятора. Данное устройство
выступает в роле конечного приёмника сообщений, обеспечивая их сбор и последующую
печать в файл. Такой подход позволяет сохранить прозрачность для пользователя,
так как гостевые операционные системы воспринимают устройство как стандартное
аппаратное обеспечение.

Основой данного механизма является общая память, разделяемая между гостевой 
операционной системой и гипервизором. Такой подход позволяет организовать
канал передачи данных без применения медленных операций ввода-вывода и без 
дополнительных копирований между буферами.

Сложность разработки заключается в том, что ради минимальных задержек
синхронизация доступа к буферу осуществляется без блокировок, благодаря 
использованию атомарных операций, что исключает риски взаимных блокировок
и обеспечивает более детерминированное время выполнения операций. Такой 
подход позволяет достичь высокой пропускной способности и минимальной задержки
при операциях записи; отсутствие блокировок в госте также может облегчить
конкуренцию за ресурсы хоста, когда на одной физической машине одновременно
работает несколько виртуальных машин.

Немаловажной задачей также является реализация механизма передачи сообщений
из среды BIOS (UEFI). На этом этапе происходит первичная инициализация 
аппаратных компонентов, и сбои приводят к невозможности дальнейшей загрузки 
системы. Отсутствие средств сбора отладочной информации существенно затрудняет
выявление причин отказов.

BIOS представляет собой интерфейс между программным обеспечением и аппаратными 
средствами, который позволяет им общаться и взаимодействовать друг с другом.
\footnote{Вонг А. Справочник по параметрам BIOS. — 2005}

UEFI (от англ. Unified Extensible Firmware Interface, единый расширяемый 
микропрограммный интерфейс) был разработан консорциумом UEFI Forum на 
основе спецификаций, изначально предоставленных корпорацией Intel как преемника
BIOS (хотя большинство образов микропрограммного обеспечения UEFI предоставляют 
поддержку для устаревших сервисов BIOS). UEFI обеспечивает поддержку современного 
оборудования, модульной архитектуры, не зависящей от архитектуры ЦПУ и драйверов, 
загрузочной сети и безопасной загрузки.
\footnote{Энциклопедия «Касперского» — UEFI}

Важным преимуществом предлагаемого решения является возможность вывода сообщений от 
гостевой операционной системы вместе с диагностическими данными от самого эмулятора.
В результате сообщения драйверов и сервисов, работающих в гостевых операционных системах,
и сообщения эмулятора печатаются в одном месте в хронологическом порядке.
Это позволяет видеть полную картину и более эффективно выявлять и исправлять ошибки.
Техподдержке и разработчикам не нужно собирать диагностическую информацию из разных
источников и пытаться сопоставить события по времени, ведь в одном файле доступны
все события, упорядоченные по времени.

Для решения проблем, описанных выше, предлагается механизм, основанный на использовании
общей памяти. На стороне эмулятора реализуется виртуальное устройство, служащее интерфейсом
для хранения метаинформации о буфере. В операционной системе разработан драйвер,
обеспечивающий взаимодействие с данным устройством и передачу диагностических данных.

Объект исследования в данной работе --- высокопроизводительный механизм передачи информации
между гостевой операционной системой и эмулятором с минимальными временными задержками.

Предмет исследования --- многопоточный кольцевой буфер без блокировок как средство передачи
диагностической информации в среде виртуализации.

Целью дипломной работы является разработка комплекса программных компонентов для организации
высокоскоростной передачи диагностической информации в среде виртуализации. В состав 
разрабатываемого решения входят: виртуальное устройство для виртуальной машины
Stratovirt
\footnote{https://www.openeuler.org/en/other/projects/stratovirt/}; драйвер для среды UEFI,
позволяющий передавать диагностические сообщения на самых ранних этапах работы эмулятора; а
также драйвер для операционной системы Windows, обеспечивающий возможность передачи диагностической
информации от приложений, сервисов и других драйверов, работающих внутри гостевой системы.

Таким образом, постановка задачи заключается в разработке данных трёх компонентов. Для достижения
поставленной цели требуется выполнить следующие этапы: \\
1. Изучить устройство и архитектуру эмулятора Stratovirt. \\
2. Разработать виртуальное устройство на стороне эмулятора. \\
3. Изучить устройство и архитектуру UEFI. \\
4. Разработать драйвер для UEFI. \\
5. Разобраться в коде некоторых драйверов для ОС Windows и изучить API, предоставляемые
операционной системой для разработки драйверов. \\
6. Реализовать драйвер для ОС Windows. \\

Методологическую основу работы составляют: анализ технической документации и литературы
о разработке драйверов для UEFI и Windows, материалы по lock-free алгоритмам и методы
отладки, применяемые в среде виртуализации. 

Постановка задачи определяет структуру дипломной работы, которая включает введение,
четыре главы и заключение.

Введение задаёт цели и задачи работы: в этой части определяются объект и предмет
исследования и раскрывается значимость.

В первой главе описывается реализация механизма передачи
диагностических данных: общая архитектура, lock-free кольцевой буфер и
интеграция компонентов.

Во второй главе рассматриваются устройство и архитектура эмулятора
Stratovirt и виртуальное устройство LogDev.

Третья глава посвящена архитектуре UEFI и разработке прошивочного
драйвера для сбора диагностики на ранних этапах загрузки.

В четвёртой главе описывается драйвер для операционной системы Windows
и интерфейсы доступа к буферу из пользовательского и ядерного режима.

Заключение содержит обобщение результатов проведённого исследования,
сравнение с другими механизмами, а также описание возможных сценариев
использования разработанного решения.

% В рамках данной работы были разработаны модуль внутри виртуальной машины 
% на базе QEMU-подобного эмулятора Stratovirt, драйвер для UEFI (EDK2), 
% а также драйвер для операционной системы Windows. Модуль внутри эмулятора 
% представлен виртуальным устройством, доступ к регистрам которого 
% осуществляется посредством memory-mapped I/O (MMIO). Устройство выполняет 
% роль приёмника сообщений, поступающих от драйверов гостевых систем, 
% осуществляет их обработку и сохраняет в файл на стороне хоста. 
% Пользовательские приложения и сервисы гостевых ОС, обращаясь к 
% соответствующим драйверам, могут передавать диагностические сообщения на 
% хост.


